\scanpage{37}\mspage{36}
% @schol
\uline{Scholie} \deleted{L'applicat\unsure{ion}} Soit $\unsure{F|\mathcal F}$ une figure. Alors

% @maponelead
a) L'application
% @mapone
\diagram{\begin{tikzcd}[ampersand replacement=\&,column sep=small]
X\arrow[r,mapsto] \& L_X=\widetilde X \& F\arrow[r,phantom,"{\unsure{\relbar\joinrel\relbar\!\shortmid|\longrightarrow}}" description] \& \unsure{\mathrm P|\mathcal P}(\unsure{L|\mathcal L})
\end{tikzcd}}
% @image
est injective, et son image \inserted{$\widetilde F$} est un ensemble de parties ``\unsure{admissible|admissibles}'' de l'ensemble $L$, i.e.
% @unions
\[
\text{\unsure{a)}}\quad\forall\widetilde X,\widetilde Y\in\widetilde F,\quad\text{on a}\quad
\widetilde X\cap\widetilde Y=\bigcup_{\substack{\widetilde Z\in\widetilde F\\ \widetilde Z\subset\widetilde X\cap\widetilde Y}}\widetilde Z
\]
% @partb
\[
\begin{aligned}
&\text{b)}\quad \forall\widetilde X\in\widetilde F,\quad\text{posant}\quad
\partial\widetilde X=\bigcup_{\substack{\widetilde Y\in\widetilde F\\ \deleted{\unsure{X|Y}}\widetilde Y\subsetneq\widetilde X}}\widetilde Y,\quad\text{et}\\
&\qquad\widetilde X^{\unsure{\circ|0}}=\widetilde X\setminus\partial\widetilde X,\quad\text{on a}\qquad\qquad \widetilde X^\circ\neq\varnothing
\end{aligned}
\]
% @margin
% slanted
\margin{\uline{\unsure{Déf}} La figure $\partial X=\{Y\in M\mid Y<X\}$ s'appelle \uline{le bord} de $X$}
% @name

On \deleted{b) L'application \unsure{$X$}} appelle $\widetilde X^\circ$ l'\uline{ombre ouverte} de $\widetilde X$ sur $L$, $\partial\widetilde X$ \deleted{\uline{l'ombre \unsure{frontière}}} \inserted{\uline{bord d'ombre}}
% @maptwolead

b) L'application
% @maptwo
\diagram{\begin{tikzcd}[ampersand replacement=\&,column sep=small]
X\arrow[r,phantom,"{\unsure{\longmapsto|\mapstochar\relbar\joinrel\relbar}}" description] \& L_X^\circ=\widetilde X^\circ \& F\arrow[r,phantom,"{\unsure{\longrightarrow|\relbar\joinrel\relbar}}" description] \& \mathrm P(L).
\end{tikzcd}}
% @partition
est injective, et \deleted{\unsure{les} de $F$} les images forment une \uline{partition} de l'ombre $L_F$ de la figure $F$, indexée par $F$.
% @order

c) \deleted{\unsure{Soit}} La relation $\leq$ sur $F$ est celle induite par la relation d'inclusion sur les ombres. Si $X\unsure{,|.}Y\in F$, alors conditions \unsure{équiv.}
% @four
\[
\left\{\begin{aligned}
&\text{(i)}\;X\leq Y\\
&\text{(ii)}\;\widetilde X\unsure{\subseteq|\subset}\widetilde Y\\
&\text{(\unsure{iii})}\;\widetilde X^\circ\subseteq\widetilde Y\\
&\text{(iv)}\;\widetilde X^\circ\cap\widetilde Y\neq\varnothing
\end{aligned}\right.
\]
